\section{Crystal Structures And CIF Import}

PyTex now supports phase construction from crystallographic structures and CIF data at the canonical data-model layer rather than only through external adapters.

\subsection{Core Rule}

If a crystallographic source defines a phase, PyTex should construct canonical primitives directly:
\begin{itemize}
  \item a lattice,
  \item a unit cell,
  \item a phase,
  \item a point-group symmetry specification derived from the source structure.
\end{itemize}

This ensures that downstream orientation, EBSD, and diffraction workflows operate on PyTex primitives rather than on source-specific structure objects.

\subsection{Current Implementation}

The present implementation uses \texttt{pymatgen} as an optional crystallographic parser and structure provider. PyTex exposes:
\begin{itemize}
  \item \texttt{Lattice.from\_pymatgen\_lattice(...)},
  \item \texttt{UnitCell.from\_pymatgen\_structure(...)},
  \item \texttt{Phase.from\_pymatgen\_structure(...)},
  \item \texttt{Phase.from\_cif(...)},
  \item \texttt{Phase.from\_cif\_string(...)}.
\end{itemize}

The imported structure is normalized into PyTex objects, including:
\begin{itemize}
  \item lattice parameters,
  \item atomic basis,
  \item reduced chemical formula,
  \item space-group symbol and number,
  \item crystal point group reduced to the supported \texttt{SymmetrySpec} surface.
\end{itemize}

\subsection{Disordered Sites}

PyTex stores one species and one occupancy per \texttt{AtomicSite}. When a source structure contains a disordered crystallographic site with multiple species, the current import path expands that source site into multiple PyTex \texttt{AtomicSite} records sharing the same fractional coordinates and carrying explicit occupancy values.

\subsection{Current Limits}

\begin{itemize}
  \item CIF support currently depends on \texttt{pymatgen} rather than on a native PyTex parser.
  \item Space-group information is retained on the phase object, but symmetry algorithms currently operate at the point-group level implemented in \texttt{SymmetrySpec}.
  \item Magnetic symmetry, superspace descriptions, and richer crystallographic metadata remain future work.
\end{itemize}

\subsection{Normative And Informative References}

\begin{thebibliography}{9}
\bibitem{iucr_cif}
International Union of Crystallography,
\textit{Crystallographic Information Framework (CIF)}.
\url{https://www.iucr.org/resources/cif}.

\bibitem{pymatgen}
S. P. Ong et al.,
\textit{Python Materials Genomics (pymatgen): A robust, open-source python library for materials analysis},
Computational Materials Science 68 (2013) 314--319.
DOI: \url{https://doi.org/10.1016/j.commatsci.2012.10.028}.
\end{thebibliography}
